I was christened Sun by my online neighbours who, finding "Suntzoogway" [sesquipedalian](https://dictionary.cambridge.org/es/diccionario/ingles/sesquipedalian), began to call me by the celestial [Anglish](https://anglish.org/wiki/Anglish) [hypocorism](https://en.wikipedia.org/wiki/Hypocorism).[^1]

Whether subconscious choice or foreshadowing by a higher storyteller,[^2] I portmanteaued Sun Tzu and Oogway without noticing the sun crowning it... but as the latter said: [there are no accidents](https://www.youtube.com/watch?v=hYAQtEs2Img&t=37s).

In Chinese, sūn 孫 means grandchild, descendant. English mirrors this with the homonym son. So sun contains both the child and the star; a compressed way of saying starchild. We're all little starchildren, little suns, as is our sun (stars all the way up). I pray we don't lose the warmth and wonder we're from.

As the son of the sun, this word weighs. But when I hear it, I am light.

*"Truly I tell you, unless you change and become like little children, you will never enter the kingdom of heaven."* — Matthew 18:3

[^1]: If this is the first sentence you read from me, don't be dismayed by the jargon! I, too, appreciate clear and concise writing. However, I'm also whimsical. I couldn't resist writing a sentence playing with the irony of doubling down on the verbosity I intend to move away from in seeking grace and beauty. In fact, Suntzugi was coined as an attempt to ease the handle... it just so happened to be richer in meaning, which is an aesthetic and philosophical pursuit of mine.
[^2]: A higher storyteller is a steller. Steller is a neologism that came to me in the feverdream of this morning's writing; the etymological child of stellar storyteller, where stella (Latin for star) carries the sense of both the heavens, the higher powers, and excellence.